% ==================== 1. Introduction ====================
\section{Introduction}

Given one image in which a scene is seen through a reflecting surface, \emph{single-image
reflection removal} (SIRR) recovers the transmission layer $T$ from the superposition
$I=T+R$ of the scene and the reflection $R$. The problem is under-determined --- two unknown
layers through one measurement --- so every method has to supply a prior. Classical approaches
assume the reflection is smoother, sparser in gradient, or of lower rank than the scene and
separate the two layers by optimisation~\cite{wan2017benchmarking,wan2022crrn}. Physics-based approaches take additional
measurements --- polarization at several angles~\cite{lei2020polarized}, or a flash/no-flash
pair --- and are accurate but need hardware that a single RGB frame does not have. Learning-based methods instead train
on paired data: synthetically composited pairs dominate (perceptual
losses~\cite{zhang2018perceptual}, cascaded refinement~\cite{li2020ibcln}, learned non-linear
coupling~\cite{hu2023dsrnet,zhao2025rdnet}, transformer and diffusion
priors~\cite{hong2024ldiffer,promptrr2024}), and real-pair corpora --- tens of thousands of
pairs captured in the wild~\cite{zhu2024rrw,openrr2025,cai2026ntire} --- have made such
captures trainable. The applications are exactly those
in which the glass cannot be taken away: photographs shot through car windows and shop
displays, museum and archive glazing, surveillance and driver-assistance cameras behind
windscreens, and product photography through protective covers.

Real-pair \emph{supervision} rests on a simple idea: capture an image without reflection and
use it as a target. The three acquisition routes --- removing the
glass, subtracting in the RAW domain, or covering the camera side with a black cloth --- were
all designed to answer \emph{how to obtain the target}, and rarely ask \emph{what properties of
the reflection layer the target silently assumes}. They also share an unstated requirement:
the camera-side cover must be \textbf{complete}. We revisit that question for a case that is
common in practice, poorly represented in existing data, and physically the least friendly to
those assumptions: \textbf{automotive tinted glass}.

\textbf{Two assumptions break on tinted glass.}

\textbf{(A2) The reflection layer is not a smooth, non-negative film.} When a foreground object
sits on the camera side --- the photographer, a handheld phone --- it blocks part of the
reflected solid angle \emph{and} part of the scene illumination. Inside the resulting
\emph{occlusion region} $\Omega$ the reflection is attenuated to about $30\%$ of its local
background, the occluder contributes a faint diffuse reflection of its own, and, most
importantly, the input drops \emph{below its own target}. Measured on our 82 confirmed pairs the two layers part company: the
reflection-free target is only $3.7\%$ darker inside $\Omega$ ($T_{\mathrm{in}}/T_{\mathrm{out}}=0.963$),
whereas the reflected content is $16.8\%$ weaker ($r_{\mathrm{in}}/r_{\mathrm{out}}=0.832$), leaving the
input $13\%$ darker than its own target there ($I/T|_{\mathrm{in}}\!\approx\!0.87$). A correction
applied at one global strength must therefore over-correct inside $\Omega$, and the effect is
visible in the errors rather than only in the model: the resulting over-correction ratio is measured at $2.7\times$. This is also why a
per-region score cannot be used to blame existing methods: their reflection field is
continuous and has no per-region gain term, so the $13\%$ input deficit has nowhere to go ---
but the region itself is not identifiable from one frame, which makes the comparison
unanswerable rather than unfavourable

\textbf{(A3) The global photometric level is not identifiable from one frame.} From
$I=\tau T+R$, the target $\tau T$ is identifiable only up to one global scale, fixed by
$\rho=\mu(\tau T)/\mu(I)$, which is not a function of the input. A pixel-wise objective
therefore degenerates toward the domain-average brightness: all 35 of our arms are
systematically $6.66$ gray levels too bright, and a closed-form brightness correction recovers
$+0.63$~dB on average ($\max$ $+1.89$~dB). The \emph{ranking} of $\rho$ within a domain is
learnable (Spearman $0.60$), its absolute value is not.

\textbf{This paper.} Rather than proposing another removal network we ask a question a
practitioner can reuse: \emph{given these two non-identifiable quantities, which supervision
actually transfers?} We answer with a benchmark, controlled arms with three seeds, a closed-form
bound, and an interventional test, and the answer is more specific than ``add an occlusion
loss''. Positively, supervised fine-tuning on our domain transfers strongly: $22.70$~dB on a
frozen 15-scene panel, $\mathbf{+2.22}$~dB over the strongest public method and better on
\textbf{every one} of the 15 scenes ($+1.69$~dB after a per-image affine, so photometry is only
about half a decibel of it) --- and the winning pool is the one containing the \emph{fewest}
occlusion regions. Negatively, the value of
an occlusion-region loss comes from its \emph{normalization} by the local reflection amplitude,
not from the region: an area-matched random region matches the true mask, a whole-image relative
error already beats it, a boundary term confined to two pixels cannot move a global score by
more than $0.036$~dB, and a mask whose labels are six times better buys no accuracy.

Our contributions are:

\begin{itemize}
  \item \textbf{A real-paired tinted-glass benchmark} of 631 pairs over 13 training scenes with
  four controlled attenuation levels, per-frame occlusion-region masks and a signed reflection
  layer, a zero-shot evaluation of seven public methods, and a controlled fine-tuning study on a
  frozen 15-scene panel (Sections~\ref{sec:data}, \ref{sec:bench}, \ref{sec:ft}).
  \item \textbf{A formalization of the occlusion region} --- piecewise smooth with sharp
  boundaries --- with an over-correction proposition, a measured decomposition of the three
  physical terms, and the non-identifiability of the global photometric scale
  (Section~\ref{sec:method}).
  \item \textbf{A diagnostic of supervision placement}: with three seeds per arm an area-matched
  random region ($19.90$~dB) matches the true mask ($19.41$~dB, $-0.49$~dB, n.s.), the
  whole-image relative error reaches $20.05$~dB, and a closed-form support-area bound explains
  why boundary-scale terms are unmeasurable by global PSNR
  (Section~\ref{sec:diag}).
  \item \textbf{Two artifacts and a calibrated negative}: a reflection mask supervised by the
  true reflection layer (IoU $0.376$ vs.\ $0.063$ for the gradient-comparison pseudo-label of
  prior work, with \emph{no} claimed downstream gain), and an interventional acceptance protocol
  under which state-of-the-art photometric conditioning scores $0.02$ against a required $0.50$
  (Section~\ref{sec:gate}, \ref{sec:photom}).
\end{itemize}

We do not claim a state-of-the-art PSNR on public benchmarks. We do claim that the two axes
above are the binding constraints on this data, that the transfer result we report is large and
reproducible, and that it is \emph{not} evidence for occlusion supervision --- so we report
every negative result that follows.

% ==================== Figure 1 (schematic + measured numbers) ====================
\begin{figure*}[t]
\centering
\begin{tikzpicture}[font=\footnotesize,>=Stealth]
  % ---- left: geometry ----
  \draw[thick] (0,0) rectangle (0.5,1.9);
  \node[align=center] at (0.25,2.15) {camera};
  \draw[thick] (2.0,-0.1) -- (2.0,2.1);
  \node at (2.0,2.35) {tinted glass};
  \draw[thick] (4.6,-0.1) -- (4.6,2.1);
  \node at (4.85,2.35) {scene};
  % unblocked reflected ray
  \draw[->,blue] (4.6,1.55) -- (2.0,1.55) -- (0.5,1.55);
  % blocked reflected ray
  \draw[->,blue] (4.6,0.55) -- (2.0,0.55) -- (1.6,0.55);
  \node[fill=white,draw,inner sep=1.5pt] (occ) at (1.15,0.55) {occluder};
  \draw[red,->] (occ.north) -- (1.15,1.05);
  \node[red] at (1.15,1.32) {blocks part of $\mathcal{A}(x)$};
  \node[red] at (1.15,1.06) {and of the light};
  \node[blue] at (4.0,0.15) {reflected path};
  % ---- right: measured ----
  \begin{scope}[xshift=6.4cm]
    \draw[->] (0,0) -- (2.5,0);
    \draw[->] (0,0) -- (0,1.9);
    \node[rotate=90,anchor=south] at (-0.32,0.95) {reflection};
    \fill[gray!35] (0.25,0) rectangle (0.95,0.57);
    \fill[gray!70] (1.25,0) rectangle (1.95,1.9);
    \node[above] at (0.6,0.57) {$\Omega$ ($30\%$)};
    \node[above] at (1.6,1.9) {outside ($100\%$)};
  \end{scope}
\end{tikzpicture}

\vspace{2pt}
\footnotesize
\begin{tabular}{@{}c@{}}
$I=\alpha T+(1-c)\beta R+cR_b+C_{\mathrm{nl}}$ \\
inside/outside $I/T=0.87$ vs $1.045$ $\Rightarrow$
\textbf{input $13\%$ below its own target} ($T_{\mathrm{in}}/T_{\mathrm{out}}=0.963$, $r_{\mathrm{in}}/r_{\mathrm{out}}=0.832$)
\end{tabular}

\caption{On tinted glass the camera-side occluder blocks a fraction $c(x)$ of the reflected
solid angle \emph{and} part of the illumination, producing a piecewise-smooth region $\Omega$
with a sharp boundary (left). Measured on 82 confirmed pairs (right): the reflection falls to
about $30\%$ of the local background and the reflected content is $16.8\%$ weaker, while the
target itself is only $3.7\%$ darker --- so the input sits $13\%$ below its own target inside
$\Omega$, which is why a single global correction strength cannot fit both regions.}
\label{fig:omega}
\end{figure*}

